\documentclass[11pt,letterpaper]{article}
\usepackage[utf8]{inputenc}
\usepackage[spanish]{babel}
\usepackage{amsmath,graphicx,float}
\usepackage[margin=2.5cm]{geometry}
\renewcommand{\baselinestretch}{1.15}
\pagenumbering{arabic}

\title{Modelos de lenguaje de n-gramas sobre caracteres:\\
       predicción en \emph{La Divina Comedia}}
\author{Equipo NN\\ Integrante A \quad Integrante B}
\date{Noviembre de 2026}

\begin{document}

\begin{titlepage}
\centering
{\Large \bfseries Modelos de lenguaje de n-gramas sobre caracteres}\\[0.6cm]
{\large Predicción en \emph{La Divina Comedia}}\\[1.2cm]
Proyecto P1 --- Probabilidad (S01)\\[0.4cm]
Equipo NN\\[0.2cm]
Integrante A: \underline{\hspace{4.5cm}}\quad Integrante B: \underline{\hspace{4.5cm}}\\[0.8cm]
Noviembre de 2026

\vfill
\textit{Borrador de trabajo: las secciones marcadas como
``Por completar: Persona B'' están pendientes de integración.}
\end{titlepage}

\begin{abstract}
Estimamos modelos de $n$-gramas de caracteres ($n=1,\dots,5$) sobre una traducción
al español de \emph{La Divina Comedia} (668\,786 caracteres normalizados). La
partición contigua 80/10/10 fija el vocabulario solo con entrenamiento
($\vert\mathcal{V}\vert=35$). El suavizado aditivo se eligió con validación; el
bloque de prueba se evaluó una sola vez. La validación seleccionó $n=5$ con
$\alpha=0.03$; en prueba el BPC fue $2.1538$ bits/carácter y la perplejidad
$4.4499$. La comparación $n=5$ frente a $n=4$ dio $\Delta=-0.1518$ bits con
intervalo bootstrap del 95\,\% $[-0.1632,-0.1410]$ (34 bloques de 2000,
10\,000 remuestras), que excluye el cero. El promedio acumulado de
$-\log_2\hat q$ se estabilizó en la figura F4 y la desigualdad de Gibbs se
cumplió con $KL=0.0016$ bits. El trabajo separa verificaciones exactas de
resultados ilustrativos y concluye que el orden superior captura dependencias
locales sin sobreajuste apreciable en este corpus.
\end{abstract}

\section{Introducción}

Predecir el siguiente carácter de un texto es el ejemplo más directo de un
modelo de lenguaje: la probabilidad de una cadena se factoriza, por la regla de
la cadena (Prop.\ 1.1.3), en el producto de condicionales de cada símbolo dado
su contexto inmediato. Cuando el contexto se limita a los $n-1$ caracteres
anteriores obtenemos una cadena de Markov de orden $n-1$ (Def.\ 1.1.6). El
proyecto pregunta cuál de esos órdenes predice mejor un texto español real,
cómo debe suavizarse y cuánta incertidumbre tiene la comparación.

Trabajamos con la traducción de Manuel Aranda y Sanjuán de \emph{La Divina
Comedia}, descargada de Internet Archive el 21 de noviembre de 2026. El texto
crudo se normalizó y se dividió en tres bloques contiguos. Todas las decisiones
de modelo ---el par $(n,\alpha)$--- se tomaron con el bloque de validación; el
bloque de prueba se reservó para una evaluación final. Este protocolo es el
que permite reportar una cifra de prueba sin haberla usado para elegir.

Los aportes de este borrador son: (i) la medición del corpus y la prueba de
cordura reproducible; (ii) los conteos MLE y la aparición de probabilidades
cero; (iii) la evaluación única en prueba con verificación de Gibbs; y (iv) un
intervalo bootstrap por bloques para la diferencia entre los dos mejores
modelos.

\section{Marco probabilístico}

Sea $\mathcal{V}$ el alfabeto finito de símbolos (Def.\ 1.1.1); aquí
$\vert\mathcal{V}\vert=35$. Escribimos $x_{1:T}$ para una cadena de caracteres
y $x_{t-n+1:t-1}$ para su contexto de longitud $n-1$. La regla de la cadena
(Prop.\ 1.1.3) da $q(x_{1:T})=\prod_t q(x_t\mid x_{1:t-1})$; el modelo de
$n$-gramas (Def.\ 1.1.6) sustituye ese contexto por los $n-1$ caracteres
previos, suponiendo Markov de orden $n-1$.

La MLE (Prop.\ 1.1.8) estima cada condicional por su frecuencia relativa:
\begin{equation}
  \hat q(y\mid h)=\frac{c(h,y)}{c(h)} ,
\end{equation}
donde $c(h,y)$ cuenta las veces que $y$ sigue al contexto $h$ en entrenamiento
y $c(h)=\sum_{y\in\mathcal{V}}c(h,y)$. Si un par no aparece, la MLE asigna
probabilidad cero. El suavizado aditivo (Prop.\ 1.1.14) corrige esto:
\begin{equation}
  \hat q_\alpha(y\mid h)=\frac{c(h,y)+\alpha}{c(h)+\alpha\vert\mathcal{V}\vert},
  \qquad \alpha=0 \text{ es la MLE}. \tag{2}
\end{equation}

La entropía de una distribución es $H(p)=-\sum_i p_i\log_2 p_i$ y satisface
$H(p)\le\log_2\vert\mathcal{V}\vert$ (Prop.\ 1.1.15), con igualdad solo si
todos los símbolos son equiprobables. La entropía cruzada o costo medio de
codificar con $q$ datos de $p$ es $H(p,q)=-\sum_i p_i\log_2 q_i$, y la
desigualdad de Gibbs (Teo.\ 1.1.27) afirma $H(p)\le H(p,q)$; la diferencia es
la divergencia de Kullback--Leibler $D_{\mathrm{KL}}(p\|q)$ (Def.\ 1.1.17).
En el modelado de caracteres se reporta
\begin{equation}
  \mathrm{BPC}=-\frac{1}{N_e}\sum_t\log_2\hat q_\alpha(x_t\mid h_t),
  \qquad \mathrm{PPL}=2^{\mathrm{BPC}}, \tag{3}
\end{equation}
donde $N_e$ es el número de caracteres evaluados (Def.\ 1.1.20). El BPC es una
entropía cruzada empírica y la PPL su exponencial en base 2.

Finalmente, el bootstrap (Def.\ 1.1.44) remuestrea con reemplazo unidades
consideradas independientes para aproximar la distribución de un estadístico.
Aquí las unidades son bloques contiguos, no caracteres, porque la cadena de
Markov introduce dependencia entre posiciones vecinas. La distribución
estacionaria $\pi$ de una matriz de transición positiva se obtiene por
iteración de potencias $\pi_0P^t$ (Teo.\ 1.1.40); ese análisis corresponde a la
Fase 5, a cargo de Persona B.

\section{Datos}

La fuente es \emph{La Divina Comedia}, Dante Alighieri, traducción de Manuel
Aranda y Sanjuán (Barcelona, 1871), descargada de
\texttt{https://archive.org/details/ladivinacomedia57303gut} el 21 de
noviembre de 2026. El archivo es texto plano en CP1252 y coincide con el
Gutenberg eBook \#57303; los detalles de procedencia están en
\texttt{datos/FUENTE.md}. El script \texttt{preparar.py} recorta el poema entre
\texttt{\_CANTO PRIMERO\_} y \texttt{\_INDICE\_}, elimina encabezados,
ilustraciones y notas, convierte a minúsculas, sustituye por un único símbolo
\# todo carácter fuera del alfabeto, colapsa espacios y fija $\mathcal{V}$ solo
sobre entrenamiento.

La tabla T1 resume la partición contigua: entrenamiento $535\,028$ caracteres,
validación $66\,879$ y prueba $66\,879$ ($668\,786$ en total), con
$\vert\mathcal{V}\vert=35$. La figura F1 muestra las frecuencias ordenadas de
unigramas. La entropía de las frecuencias de entrenamiento es
$H(\hat p)=4.1992$ bits, mientras que el máximo posible es
$\log_2 35=5.1293$ bits. La diferencia de $0.9301$ bits mide cuánto se aleja el
texto de un sorteo uniforme sobre el alfabeto: el español usa de forma muy
desigual las letras y, sobre todo, el espacio. En términos relativos la
entropía observada es aproximadamente $81.8\,\%$ del máximo, por lo que el texto está lejos
de ser aleatorio, aunque aún conserva casi un bit de incertidumbre por símbolo
bajo el modelo unigramas.

\section{Métodos}

\textbf{[Por completar: Persona B]} Debe describir los conteos, el suavizado,
la selección de $\alpha$ con F2, las métricas, la simulación y el bootstrap,
con suficiente detalle para repetirlos.

\section{Resultados}

\textbf{[Por completar: Persona B]} Debe citar y comentar T2--T4 y F2--F5.
Los archivos T2, T3, T4, F3 y F4 ya están generados; F2 y F5 dependen de la
Persona B.

\section{Verificaciones}

\textbf{[Por completar: Persona B]} Debe reportar la salida de
\texttt{pruebas.py}: los seis valores de \emph{abracadabra}, las sumas
condicionales, Gibbs y la convergencia de $\pi$.

\section{Discusión}

\subsection{Distancia al azar (Fase 1)}

La comparación $H(\hat p)=4.1992$ frente a $\log_2 35=5.1293$ bits no significa
que el texto ``contenga'' menos información que una cadena aleatoria: significa
que la distribución empírica de símbolos es más concentrada. Codificando con un
modelo uniforme se gastarían $5.1293$ bits por carácter; usando las frecuencias
reales se gastan $4.1992$, un ahorro de $0.9301$ bits por símbolo. El modelo
deberá explotar además la dependencia entre caracteres, que es lo que miden los
BPC de T3.

\subsection{Conteos MLE y probabilidades cero (Fase 2)}

La tabla T2 muestra que el número de contextos distintos crece de $1$ ($n=1$) a
$23\,672$ ($n=5$), mientras que la fracción de caracteres de validación con
probabilidad cero bajo la MLE sube de $0$ a $0.0391$. El BPC de la MLE en
validación es infinito en cuanto aparece un solo cero porque el BPC promedia
$-\log_2\hat q$ y $-\log_2 0=+\infty$; un solo evento imposible domina la media
por completo. Esto no es un defecto de la métrica sino una consecuencia de
haber asignado probabilidad cero a un evento que sí ocurrió.

Los ceros crecen con $n$ porque el espacio de contextos posibles crece como
$\vert\mathcal{V}\vert^{n-1}$ y las $535\,028$ posiciones de entrenamiento no
pueden cubrirlo. Aunque los contextos realmente observados son muchos, las
transiciones rara vez se repiten: para $n=5$ hay $66\,833$ pares observados
repartidos entre $23\,672$ contextos, es decir, menos de tres sucesos por
contexto. La perturbación de $0.01$ de probabilidad en un contexto frecuente
confirmó la Prop.\ 1.1.8: en los cinco órdenes la MLE tuvo mayor
log-verosimilitud que las dos alternativas perturbadas.

\subsection{Elección de orden y evaluación en prueba (Fase 4)}

La validación eligió $n=5$ ($\alpha=0.03$) porque dio el menor BPC de
validación, $2.1838$ bits, frente a $4.2000$, $3.2689$, $2.7432$ y $2.3362$ de
los órdenes $1$ a $4$. La prueba, evaluada una sola vez, confirma el mismo
orden: los BPC de T3 decrecen monótonamente hasta $2.1538$ bits en $n=5$ y la
perplejidad correspondiente es $4.4499$. Por tanto, la prueba no contradice la
selección de validación.

En cambio, el BPC de entrenamiento sigue bajando hasta $1.8482$ bits en $n=5$.
La explicación es la misma que en cualquier modelo con muchos parámetros: al
crecer $n$, cada condicional se estima con menos datos y se ajusta a
idiosincrasias del entrenamiento. Validación y prueba pagan ese coste y por eso
su curva se aplana. La figura F3 muestra ambas curvas en la misma escala y la
figura F4 muestra que el promedio acumulado de $-\log_2\hat q$ se estabiliza,
como predice la ley de los grandes números. La desigualdad de Gibbs también se
cumplió con el modelo unigramas: $H(\hat p_{\mathrm{prueba}})=4.1909\le
H(\hat p_{\mathrm{prueba}},\hat q)=4.1925$, con $KL=0.0016$ bits.

\subsection{Incertidumbre de la comparación (Fase 6)}

La tabla T4 reporta $\Delta=-0.1518$ bits a favor de $n=5$ frente a $n=4$, con
intervalo bootstrap del 95\,\% $[-0.1632,-0.1410]$ basado en $34$ bloques
contiguos de $2000$ caracteres y $10\,000$ remuestras. El intervalo excluye el
cero, así que la mejora observada no es compatible con una diferencia nula bajo
este esquema de remuestreo. Esto concuerda con la pequeña separación entre los
BPC de prueba de ambos modelos ($2.1538$ frente a $2.3049$ bits).

Se remuestrean bloques y no caracteres individuales porque los caracteres no
son intercambiables: la cadena de Markov genera dependencia local y un
bootstrap de caracteres destruiría precisamente la estructura que el modelo
intenta aprovechar, subestimando la variabilidad. Los bloques contiguos de
$2000$ caracteres conservan esa dependencia a corto plazo y son suficientemente
largos para que la media de bloque sea estable. El último bloque tiene menos
caracteres que los demás, pero su contribución se promedia por su longitud, no
se descarta.

\subsection{Pendiente de la Persona B}

\textbf{[Por completar: Persona B]} Debe responder las preguntas en cursiva de
las fases 3 y 5: por qué cambia $\alpha$ con $n$ y qué ocurre con $\alpha=1$ en
contextos raros; y por qué la distribución estacionaria se parece a las
frecuencias de unigramas y qué capturan o no las muestras generadas.

\section{Conclusiones}

En este corpus el orden $n=5$ con suavizado $\alpha=0.03$ obtuvo el menor BPC
de validación y también el menor BPC de prueba, $2.1538$ bits por carácter
($\mathrm{PPL}=4.4499$). La ventaja sobre $n=4$, de $0.1518$ bits, tiene un
intervalo bootstrap del 95\,\% que excluye el cero. La prueba de
\emph{abracadabra}, las sumas condicionales y la desigualdad de Gibbs pasaron,
de modo que las cifras principales son reproducibles. Como limitación, el
modelo sigue siendo una cadena de Markov sobre caracteres: captura ortografía y
ritmo local, pero no semántica ni dependencias largas.

\section*{Referencias}

Dante Alighieri. \emph{La Divina Comedia}, traducción de Manuel Aranda y
Sanjuán. Barcelona: Empresa Editorial La Ilustración, 1871. Copia en Internet
Archive: \texttt{https://archive.org/details/ladivinacomedia57303gut}.

Efron, B. y Tibshirani, R. J. (1993). \emph{An Introduction to the Bootstrap}.
Chapman \& Hall.

Cover, T. M. y Thomas, J. A. (2006). \emph{Elements of Information Theory},
2.a edición. Wiley.

Marcano, J. \emph{Apuntes de Probabilidad}, sección S01: modelos de lenguaje,
entropía y bootstrap. Material del curso.

Shannon, C. E. (1948). A Mathematical Theory of Communication.
\emph{Bell System Technical Journal}, 27, 379--423.

\appendix

\section{Deducción de la MLE}

Para un contexto $h$ fijo queremos maximizar
$\ell(\mathbf{q})=\sum_{y\in\mathcal{V}}c(h,y)\log q(y)$ sujeto a
$\sum_{y\in\mathcal{V}}q(y)=1$. Con un multiplicador de Lagrange $\lambda$,
\[
  \frac{\partial}{\partial q(y)}
  \left[\sum_{y}c(h,y)\log q(y)+\lambda\left(1-\sum_y q(y)\right)\right]
  =\frac{c(h,y)}{q(y)}-\lambda=0 .
\]
Luego $q(y)=c(h,y)/\lambda$. Imponiendo la restricción,
$\sum_y c(h,y)/\lambda=1$, de donde $\lambda=c(h)=\sum_y c(h,y)$ y por tanto
$\hat q(y\mid h)=c(h,y)/c(h)$, que es la fórmula usada en el código. La
solución es única porque $\ell$ es cóncava en el simplex. Con suavizado
aditivo se maximiza la misma verosimilitud con un previo de Dirichlet de
parámetros $\alpha$; la media posterior es la ecuación (2).

\section{Reproducibilidad}

Entorno local: Python 3.14.0, NumPy 2.2.6, Matplotlib 3.10.7. Comandos:
\begin{verbatim}
python -m pip install -r requirements.txt
python codigo/preparar.py
python codigo/pruebas.py
python codigo/experimentos.py A
\end{verbatim}
El orden de los modelos es $n=1,\dots,5$; los valores de suavizado probados en
validación son $0.001,0.003,0.01,0.03,0.1,0.3,1$. Los $\alpha$ elegidos por
validación quedan en \texttt{resultados/alpha_por_n.json}, archivo intermedio
que se regenera y no se entrega. La semilla del bootstrap es $20261122$ y el
número de remuestras es $10\,000$. La de la simulación de la Fase 5 la reportará
Persona B. El bloque de prueba se evaluó una sola vez en
\texttt{evaluar\_prueba\_una\_vez}; una segunda ejecución de la Fase 4 imprime
que la prueba no se reevalúa. Solo el flag \texttt{--forzar-prueba} permite
regenerarla deliberadamente, y no se usó para elegir $n$ ni $\alpha$.

\section{Muestras generadas y declaración de IA}

\textbf{Muestras (Persona B):} tres cadenas de 300 caracteres de los modelos
$n=2,3,5$, con la semilla que se reporte en el Anexo B.
\textbf{[Por completar: Persona B]}

\paragraph{Declaración de uso de IA.}
Se usó el asistente OpenCode (modelo Big Pickle) para planificar el trabajo,
redactar el esqueleto de los scripts, revisar las fórmulas y producir borradores
de texto. Las verificaciones propias fueron: ejecutar \texttt{pruebas.py} y
cotejar los seis valores de \emph{abracadabra}; comprobar que las condicionales
suman $1$ con error máximo $1.11\times10^{-16}$; ejecutar dos veces las fases A
y confirmar que la segunda no reescribe \texttt{metricas\_prueba.csv}; e
inspeccionar T2, T3 y T4 para que cada cifra del reporte provenga de un archivo
generado por el código. Ni la fuente ni las cifras se tomaron de la salida del
asistente sin esta comprobación.

\section*{Tabla de contribución}

\begin{center}
\begin{tabular}{lll}
\hline
Integrante & Código y resultados & Reporte \\
\hline
A & Fase 1, 2, 4 y 6; T1--T4, F1, F3, F4 & R1, R3 (parte A), R4, Anexos A--C \\
B & Fase 3, 5 y 7; F2, F5 y muestras & R2, R3 (parte B), R5 y zip \\
Ambos & Revisión de \texttt{modelos.py} y \texttt{pruebas.py} & Revisión cruzada \\
\hline
\end{tabular}
\end{center}

\end{document}
